\documentclass[../../../include/open-logic-section]{subfiles}

\begin{document}

\olfileid{sth}{replacement}{limofsize}
\olsection{Watesan Ukuran}

Mbokmenawa upaya kang paling lumrah kanggo
menehi dhasar ``intrinsik'' marang Panggantian
nganggo gagasan iki:
\begin{enumerate}
	\item[] \limofsize. Barang-barang apa wae
	mbentuk himpunan, anggere cacahe ora kakehan.
\end{enumerate}
Prinsip iki katon langsung ndhukung Panggantian:
citra saka sawijining himpunan minangka citra
surjektif saka himpunan asale. Luwih persis,
upamane kita mbentuk himpunan
$\funimage{\tau}{A} = \Setabs{\tau(x)}{x \in A}$
kanthi Panggantian. Nanging pratelan
$\cardle{\funimage{\tau}{A}}{A}$, yen
``ora luwih gedhe'' ditegesi liwat injeksi,
ora ndherek tanpa Aksioma Pilihan. Yen prinsip
``ora kakehan'' uga lumaku kanggo citra surjektif,
mula yen !!{element}s saka $A$ ora kakehan kanggo
mbentuk himpunan, citrane uga ora kakehan
kanggo mbentuk $\funimage{\tau}{A}$.

Kangelan kang cetha nalika nggunakake
\limofsize{} kanggo menehi dhasar marang
Panggantian yaiku kita \emph{durung} netepake
prinsip kaya \limofsize. Kajaba iku, nalika
kita nyritakake konsepsi himpunan kang
kumulatif lan iteratif ing
\crefrange{sth:story::chap}{sth:z::chap}, ora
ana bab kang \emph{menehi pratandha} tumuju
\limofsize. Mula Boolos tau nulis:
``Mbokmenawa wong bisa nyimpulake yen paling
ora ana rong gagasan `ing mburine' teori
himpunan'' \citeyearpar[p.~19]{Boolos1989}.
Ing sisih siji, gagasan babagan konsepsi
himpunan kang kumulatif lan iteratif dimaksudake
kanggo ndhukung~$\Z$. Ing sisih liyane,
\limofsize{} dimaksudake kanggo ndhukung
Panggantian.

Nanging masalahe ora mung amarga nganti saiki
kita \emph{durung ngrembug} \limofsize. Masalah
kang luwih dhasar yaiku \limofsize{} (miturut
rumusan ing ndhuwur) katon ora selaras karo
gagasan himpunan kang kumulatif lan iteratif.
Prinsip iki ora nyebutake pambentukan himpunan
lumantar \emph{tataran}.

Iki ora nggumunake yen ndeleng sejarah ``rong
gagasan'' mau (yaiku konsepsi himpunan kang
kumulatif lan iteratif, lan \limofsize).
Pungkasane, ``rong gagasan'' iki minangka
rong cara kang rada beda kanggo nyingkiri
paradoks teori himpunan. Konsepsi kumulatif
lan iteratif nyingkiri Paradoks Russell kanthi
kandha, kanthi garis gedhe:
\emph{kita ora kena ngarepake anane himpunan
Russell, amarga ora bakal ``kabentuk'' ing
tataran apa wae}. Kosokbaline, \limofsize{}
dimaksudake kanggo nyingkirake himpunan
Russell kanthi kandha:
\emph{kita ora kena ngarepake anane himpunan
Russell, amarga ukurane bakal kegedhen}.

Yen mangkono, ayo kandha kanthi terang:
minangka jawaban tumrap paradoks, \limofsize{}
mbutuhake dhasar kang \emph{adoh} luwih kuwat.
Coba, upamane, versi Paradoks Russell iki:
\emph{ora ana asu pug kang ngambu persis asu-asu
pug kang ora ngambu awake dhewe} (delengen
\olref[sth][story][rus]{sec}). Yen ana kang
takon ``apa sebabe asu pug kaya mangkono ora
ana?'', jawaban yen asu iku kudu ngambu kakehan
asu pug ora trep. Dadi apa sebabe katrangan
intuitif kanggo ora anane himpunan Russell
bakal trep yen mung kandha yen himpunan iku
``kegedhen'' kanggo ana?

Dadi, bisa dimaklumi yen wong rada bingung
marang dhasar ``intuitif'' kanggo \limofsize.

\end{document}
